%% Abstract and key points (written for the white paper; all claims are taken from the body)
\section*{Abstract}
\addcontentsline{toc}{section}{Abstract}
\markboth{Abstract}{}

Sixth-generation networks are expected to be managed at the timescales on which they change, by
intelligence placed throughout the infrastructure rather than concentrated above it. In UNITY-6G that
intelligence is contributed by many partners, addresses different domains---the core, the open radio
access network, non-terrestrial segments, wireless transport, non-3GPP access, the AI/ML domain,
vertical applications and digital twins---and is realized with techniques ranging from rules and
solvers through trained models to language-model reasoning. Yet it acts on shared infrastructure, often
on the same resources, toward objectives that pull against each other.

This white paper specifies the \emph{AI-native agentic framework} under which such intelligence is
hosted. The unit of hosting is the \emph{agent}: a composition of four capabilities---the Monitoring
System, the Analytics Engine, the Decision Engine and the Actuator---instantiated against a declared
objective, described by a signed descriptor, and holding only the authority the platform grants it.
The framework treats authority to change the network as a shared resource: granted at admission,
bounded by a declared scope and action space, held in a registry, mediated by a conflict fabric before
any intent reaches the network, verified against independently observed state, and revocable at any
time. Its contracts are stated precisely enough to be implemented without further interpretation:
every declaration has a schema, every check a point at which it runs, and every obligation a party
that carries it.

\begin{keybox}[title={Key points}]
\begin{itemize}
  \item \textbf{Authority is a platform resource.} An agent never grants itself authority and never
        approves its own decision; authority is allocated at admission, bounded and revocable.
  \item \textbf{Four capabilities, each defined by what it never produces.} The Monitoring System
        produces evidence, never conclusions; the Analytics Engine insight, never action; the Decision
        Engine intent, never a controller command; the Actuator executes, verifies and reports one of
        four typed outcomes.
  \item \textbf{One contract per agent.} The agent descriptor declares what an agent is, what it may
        touch, how fast it runs, where it may be placed, how it fails and who answers for it.
  \item \textbf{Mediation before actuation.} Intents that could conflict are resolved synchronously and
        fail-closed by the conflict fabric, off the per-decision fast path.
  \item \textbf{Verified, honest assurance.} An acknowledgement is not proof of change; outcomes that
        cannot be verified are reported as unverified, and assurance is built from platform events.
  \item \textbf{Technique-neutral and standards-bound.} The same contracts hold whether a capability is
        realized by rules, solvers, trained models or language models, and the framework binds to the
        O-RAN latency classes and the 3GPP management vocabulary.
\end{itemize}
\end{keybox}

\subsection*{How this paper is organized}
Section~\ref{introduction} introduces the problem and the agent; Section~\ref{lineage-and-state-of-the-art}
traces the lineage of the four-capability pattern and states what is new. Section~\ref{the-ai-native-agent-framework}
defines the framework itself---terms, principles, taxonomy, the agent descriptor, the execution model,
latency classes and the authority bound stack. Sections~\ref{the-monitoring-system}--\ref{the-actuator}
specify each of the four capabilities, and Sections~\ref{interconnection}--\ref{integration} cover
interconnection, lifecycle, conflict resolution and governance, and integration.
Section~\ref{key-performance-indicators-and-value-indicators} relates the framework to the project's
key performance and value indicators. Normative requirements appear as numbered requirement cards;
Annex~\ref{annex-a-json-schemas-of-the-agent-descriptor-and-related-documents} gives the JSON Schemas
of the descriptor and Annex~\ref{annex-b-initial-content-of-the-type-catalog} the initial content of the
type catalog.
